% ============================================================
% PROOF: The Coupling Polynomial and the Torsion-Free Split
% Supplementary material for Letter 3
% ============================================================

\documentclass[11pt]{article}
\usepackage{amsmath,amssymb}
\usepackage[margin=1in]{geometry}

\title{The Coupling Polynomial Generates the 40/24 Split}
\author{Stephen Nelson\\Field Dynamics Research}
\date{}

\begin{document}
\maketitle

\begin{abstract}
The torsion-free condition
$\Gamma^\lambda_{\mu\nu} = \Gamma^\lambda_{\nu\mu}$
retains 40 of the 64 affine connection components and
discards~24.  We show that this split is generated by the
coupling polynomial $f(K,s) = 1 + sK + K^2$ introduced in
the companion letter~\cite{Nelson2026b}, evaluated at its two
opposite states: the retained count equals
$(f(K,{+}1) - 1)/2$ per face, and the discarded count equals
$(f(K,{-}1) - 1)/2$ per face.  The derivation uses only the
polynomial and the definition of symmetric and antisymmetric
index pairs.  The physical identification of the retained
projection with TetA and the discarded projection with TetB
is supplied by the companion
letters~\cite{Nelson2026a,Nelson2026b}.  No new postulates
are introduced.
\end{abstract}

% ============================================================
\section{Setup}
% ============================================================

The affine connection on a $K$-dimensional frame has three
indices $\Gamma^\lambda_{\mu\nu}$, each running from 0 to
$K - 1$.  The total number of components is $K^3$.  At
$K = 4$, this gives $64$.

The torsion-free condition forces symmetry of the lower pair:
\begin{equation}
  \Gamma^\lambda_{\mu\nu} = \Gamma^\lambda_{\nu\mu}.
  \label{eq:torsion_free}
\end{equation}
This is the standard constraint adopted in Riemannian
geometry and inherited by general relativity.

The companion letters established:
\begin{itemize}
  \item Letter~1~\cite{Nelson2026a}: the tetrad equation produces
    a stellated octahedron geometry with structural constants
    $K = 4$, $D = 3$, selected by the uniqueness equation
    $2(D{+}1) = 2^D$.
  \item Letter~2~\cite{Nelson2026b}: the tetrad equation
    $g_{\mu\nu} = \eta_{ab}\,e^a{}_\mu\,e^b{}_\nu$ contains
    three structural ingredients at ranks 0, 1, and~2 on the
    frame space, with graded dimensions $K^0 = 1$, $K^1 = K$,
    and $K^2$.  The tetrad equation is quadratic in the frame
    field and admits two roots related by spatial parity
    $e^i{}_\mu \to -e^i{}_\mu$ ($i = 1,2,3$).  The metric is invariant under
    this parity (both roots produce the same $g_{\mu\nu}$), so
    the metric cannot distinguish the two roots.  Under parity,
    a rank-$r$ tensor transforms with sign $(-1)^r$.  Ranks~0
    and~2 are parity-even; rank~1 is parity-odd.  The rank-1
    vielbein is therefore the only ingredient whose sign
    changes when the two roots are exchanged.  This sends
    $s \to -s$, identifying TetA with $s = +1$ and TetB with
    $s = -1$.  The coupling polynomial is:
    \begin{equation}
      f(K,s) = 1 + sK + K^2.
      \label{eq:poly}
    \end{equation}
\end{itemize}

% ============================================================
\section{The split}
% ============================================================

\subsection*{Step 1: The $K \times K$ matrix}

For each value of the upper index $\lambda$, the lower pair
$(\mu, \nu)$ forms a $K \times K$ matrix with $K^2 = 16$
entries.

\subsection*{Step 2: Diagonal and off-diagonal}

The 16 entries split into:
\begin{itemize}
  \item $K = 4$ diagonal entries ($\mu = \nu$).  These have no
    mirror partner: the pair $(\mu, \mu)$ reversed is itself.
    They are inherently symmetric.
  \item $K^2 - K = 12$ off-diagonal entries ($\mu \neq \nu$).
    These form $\tfrac{1}{2}(K^2 - K) = 6$ mirror pairs:
    each $(\mu, \nu)$ is paired with $(\nu, \mu)$.
\end{itemize}

\subsection*{Step 3: The torsion-free condition splits each
mirror pair}

For each of the 6 mirror pairs, the torsion-free
condition~(\ref{eq:torsion_free}) forces
$\Gamma^\lambda_{\mu\nu} = \Gamma^\lambda_{\nu\mu}$.  This
retains the symmetric combination and discards the
antisymmetric combination:
\begin{align}
  \text{Retained:} \quad
    & \tfrac{1}{2}(\Gamma^\lambda_{\mu\nu}
    + \Gamma^\lambda_{\nu\mu}),
    \label{eq:sym} \\[4pt]
  \text{Discarded:} \quad
    & \tfrac{1}{2}(\Gamma^\lambda_{\mu\nu}
    - \Gamma^\lambda_{\nu\mu}).
    \label{eq:antisym}
\end{align}
From each pair, one combination is kept and one is discarded.

\subsection*{Step 4: The count per face}

\begin{align}
  \text{Kept per face} &= K + 6
    = K + \frac{K^2 - K}{2}
    = \frac{K^2 + K}{2}
    = \frac{16 + 4}{2} = 10,
    \label{eq:kept} \\[6pt]
  \text{Discarded per face} &= 6
    = \frac{K^2 - K}{2}
    = \frac{16 - 4}{2} = 6.
    \label{eq:disc}
\end{align}

\subsection*{Step 5: Total across all faces}

Multiplying by $K = 4$ observation faces:
\begin{align}
  \text{Total kept} &= K \times \frac{K^2 + K}{2}
    = \frac{K^3 + K^2}{2}
    = \frac{64 + 16}{2} = 40,
    \label{eq:total_kept} \\[6pt]
  \text{Total discarded} &= K \times \frac{K^2 - K}{2}
    = \frac{K^3 - K^2}{2}
    = \frac{64 - 16}{2} = 24.
    \label{eq:total_disc}
\end{align}
Check: $40 + 24 = 64 = K^3$.\quad$\checkmark$

% ============================================================
\section{Connection to the coupling polynomial}
% ============================================================

\subsection*{Step 6: The polynomial identity}

The coupling polynomial~(\ref{eq:poly}) evaluates to:
\begin{align}
  f(K, +1) &= 1 + K + K^2 = 21,
    \label{eq:f_plus} \\
  f(K, -1) &= 1 - K + K^2 = 13.
    \label{eq:f_minus}
\end{align}

Subtract the field-origin term ($K^0 = 1$) from each.  The
subtraction is not an arithmetic convenience.  The $K^0 = 1$
term of the polynomial counts the rank-0 structural ingredient
of the tetrad equation: the scalar closure condition
$\det(e) \neq 0$~\cite{Nelson2026b}.  In the connection
matrix, this maps to the diagonal entries ($\mu = \nu$), which
have no mirror partner: the pair $(\mu, \mu)$ reversed is
itself.  These entries are invariant under index exchange and
therefore belong to neither the symmetric nor the antisymmetric
sector.  They are shared (root-blind).  Removing them before
halving separates the root-invariant structure from the
root-sensitive structure that the torsion-free condition
splits:
\begin{align}
  f(K, +1) - 1 &= K + K^2 = K^2 + K = 20,
    \label{eq:fminus1_plus} \\
  f(K, -1) - 1 &= -K + K^2 = K^2 - K = 12.
    \label{eq:fminus1_minus}
\end{align}

Halve each (the mirror-pair projection):
\begin{align}
  \frac{f(K, +1) - 1}{2} &= \frac{K^2 + K}{2} = 10,
    \label{eq:half_plus} \\[6pt]
  \frac{f(K, -1) - 1}{2} &= \frac{K^2 - K}{2} = 6.
    \label{eq:half_minus}
\end{align}

Comparing~(\ref{eq:half_plus}) with~(\ref{eq:kept}) and
(\ref{eq:half_minus}) with~(\ref{eq:disc}):
\begin{equation}
  \boxed{%
  \text{Kept per face}
    = \frac{f(K,{+}1) - 1}{2}, \qquad
  \text{Discarded per face}
    = \frac{f(K,{-}1) - 1}{2}.}
  \label{eq:main}
\end{equation}

\subsection*{Step 7: Interpretation}

The expressions in~(\ref{eq:main}) are identical.  The
combinatorial formula for the symmetric/antisymmetric split
of a $K \times K$ matrix is the coupling polynomial evaluated
at its two opposite states ($s = \pm 1$), with the
field-origin term removed and the mirror-pair degeneracy
halved.

The $\pm K$ that separates 10 from 6 is the $sK$ term of the
polynomial: the parity-odd rank-1 contribution.  In the
polynomial, $sK$ counts the $K = 4$ vielbein legs, whose sign
flips under the parity that exchanges the two
roots~\cite{Nelson2026b}.  In the connection matrix, the
off-diagonal entries ($\mu \neq \nu$) are the only entries
affected by index exchange $(\mu, \nu) \leftrightarrow
(\nu, \mu)$.  The $sK$ term and the off-diagonal variance
measure the same thing: the part of the structure that is
sensitive to the ordering of two directions.  The symmetric
off-diagonal combinations (upper triangle, $+K$) are retained;
the antisymmetric combinations (lower triangle, $-K$) are
discarded.  This is why $\pm K$ splits the count: it is the
parity-sensitive sector of both the polynomial and the matrix.

The torsion-free condition therefore:
\begin{enumerate}
  \item retains the $s = +1$ projection of the coupling
    polynomial (the forward reading of each mirror pair), and
  \item discards the $s = -1$ projection (the reverse reading).
\end{enumerate}

Because the companion letters proved that spatial parity sends
$s \to -s$~\cite{Nelson2026a,Nelson2026b}, identifying TetA with $s = +1$ and TetB with
$s = -1$, the retained projection is the TetA contribution
and the discarded projection is TetB.  The 24~components
removed by the torsion-free condition are not abstract
antisymmetric index pairs.  They are the second root of the
tetrad equation, carrying a distinct field fraction
$1 - \mu(x)$ at every radius~\cite{Nelson2026b}.

% ============================================================
\section{The difference}
% ============================================================

The surplus of kept over discarded, per face, is:
\begin{equation}
  \frac{f(K,{+}1) - 1}{2} - \frac{f(K,{-}1) - 1}{2}
  = \frac{f(K,{+}1) - f(K,{-}1)}{2}
  = \frac{\Delta f}{2}
  = \frac{2K}{2}
  = K = 4.
  \label{eq:diff_per_face}
\end{equation}

Across all $K$ faces:
\begin{equation}
  K \times K = K^2 = 16.
  \label{eq:diff_total}
\end{equation}

The total surplus of kept over discarded is $K^2 = 16$: the
bilinear coupling term of the polynomial.  The asymmetry
between 40 and~24 is the $K^2$ term itself.

% ============================================================
\section{Field closure}
% ============================================================

Restoring the discarded projection returns
$K^2 = 16$ components per face:
\begin{equation}
  \frac{f(K,{+}1) - 1}{2} + \frac{f(K,{-}1) - 1}{2}
  = \frac{(K^2 + K) + (K^2 - K)}{2}
  = K^2 = 16.
  \label{eq:restored}
\end{equation}

Across all $K$ faces: $K \times K^2 = K^3 = 64$.

The full affine connection is recovered.  Both projections of
the coupling polynomial are active.  The torsion-free
condition is relaxed.  This is field closure.

% ============================================================
\section{Summary}
% ============================================================

\begin{center}
\renewcommand{\arraystretch}{1.3}
\begin{tabular}{lccc}
\hline
  & Per face & Formula & Total ($\times K$) \\
\hline
Kept (Christoffel)
  & 10 & $(f(K,{+}1){-}1)/2$ & 40 \\
Discarded (torsion)
  & 6 & $(f(K,{-}1){-}1)/2$ & 24 \\
Surplus (kept $-$ disc.)
  & $K = 4$ & $\Delta f / 2$ & $K^2 = 16$ \\
Full connection
  & $K^2 = 16$ & $K^2$ & $K^3 = 64$ \\
\hline
\end{tabular}
\end{center}

\medskip

The torsion-free condition is the $s = +1$ projection of the
coupling polynomial.  The torsion it removes is the $s = -1$
projection.  The polynomial generates both from one equation
and one sign flip.  Restoring both projections returns the
full $K^3 = 64$ component connection.

\begin{thebibliography}{9}

\bibitem{Nelson2026a}
S.~Nelson,
\textit{The Coordinate-Space Geometry of the Tetrad Frame},
Class.\ Quantum Grav.\ (2026), submitted (CQG-115428).

\bibitem{Nelson2026b}
S.~Nelson,
\textit{The Other Half of Gravity: Torsion and the
Two-Chamber Structure of the Tetrad Frame},
Class.\ Quantum Grav.\ (2026), submitted (CQG-115447).

\end{thebibliography}

\end{document}
